\documentclass{exam}

\newif\ifA
\newif\ifkey
\newif\ifprintselected   % required by exam_config's \ansbox

%%%%%%%%%%%%%%%%%%%%%%%%%%%%%%%%%%%%%%%%%%%%%%%%%%%%%%
%%% SET THESE IF VARIABLES
%%%%%%%%%%%%%%%%%%%%%%%%%%%%%%%%%%%%%%%%%%%%%%%%%%%%%%
% Lab name -- labs are named by the lesson they belong to, never by a
% sequential lab number (Neil, 2026-10-12): this is the L7 Lab.
\def\labLesson{7}

% is this the key or not?
\ifdefined\iskey
	\ifnum\iskey=1
		\keytrue
	\else
		\keyfalse
	\fi
\else
	%% Pick one manually
	\keyfalse % if this is NOT the key
%	\keytrue     % if this IS the key
\fi

% set up the header and footer - change for every term
\ifdefined\thisTermIn
	\def\thisterm{\thisTermIn}
\else
	\def\thisterm{Spring 2027}
\fi
%
\def\thisclass{ECE 448}

%%%%%%%%%%%%%%%%%%%%%%%%%%%%%%%%%%%%%%%%%%%%%%%%%%%%%%
%%% END SET THESE IF VARIABLES
%%%%%%%%%%%%%%%%%%%%%%%%%%%%%%%%%%%%%%%%%%%%%%%%%%%%%%
% packages and shortcuts
\input{myPackages_gr}
\usepackage[hidelinks, linkcolor=red]{hyperref}
\input{myShortcuts}
\setlength{\parindent}{0pt}
\setlength{\parskip}{6pt}

%% exam class customization
\input{exam_config}
% ECE 448 document fonts.
%
% Body text is Barlow (SIL OFL; the faces and the license are in
% latex/fonts/Barlow/). Barlow is not in TeX Live, so it is loaded as an
% OpenType face through fontspec -- which means the document must be built
% with lualatex, not pdflatex. build_practice.sh and build_lab.sh both do.
%
% Compiling with pdflatex anyway is not an error: the guard below falls back
% to the LaTeX default so the source still compiles under any engine and can
% be syntax-checked. The resulting PDF just is not the one we ship.
%
% Paths are relative to latex/, which is where both build scripts run
% (each does cd "$(dirname "$0")"). Loading by path rather than by installed
% family name means no system font install and no fontconfig, so a fresh
% clone builds a byte-comparable PDF anywhere.
%
% MATH IS DELIBERATELY LEFT IN COMPUTER MODERN. Barlow has no math companion,
% and CM math under sans body text is already the house pairing: the decks set
% words in Source Sans Pro and let MathJax render symbols in its CM-derived
% face. Moving math to a sans as well would mean unicode-math plus a sans math
% font (Fira Math being the plausible one) and re-checking every equation in
% the course for symbol coverage. That is a separate decision, not a side
% effect of this one.

\usepackage{iftex}
\ifLuaTeX
	\usepackage{fontspec}
	\defaultfontfeatures{Ligatures=TeX}
	\setmainfont{Barlow}[
		Path           = fonts/Barlow/,
		Extension      = .ttf,
		UprightFont    = *-Regular,
		ItalicFont     = *-Italic,
		BoldFont       = *-Bold,
		BoldItalicFont = *-BoldItalic,
	]
	\setsansfont{Barlow}[
		Path           = fonts/Barlow/,
		Extension      = .ttf,
		UprightFont    = *-Regular,
		ItalicFont     = *-Italic,
		BoldFont       = *-SemiBold,
		BoldItalicFont = *-BoldItalic,
	]
	% microtype earns its place here specifically because of Barlow. Its
	% glyphs are narrower than Computer Modern's, so justified prose sets
	% tighter and TeX has to squeeze interword space to fit; two paragraphs
	% in the L02/L03 sets went overfull on the first Barlow build. Protrusion
	% and font expansion absorb that -- both are lualatex features, so this
	% belongs inside the guard with everything else.
	\usepackage{microtype}

	% Mono stays Latin Modern: Barlow ships no monospaced face, and verbatim
	% and code need one that actually is fixed-pitch.
\fi
   % Barlow body text (needs lualatex)
\printselectedfalse      % full worked solutions in the key, not answers-only
\CorrectChoiceEmphasis{\color{red}}
\SolutionEmphasis{\color{red}}
\renewcommand\questionlabel{\textbf{\thequestion.}}
\renewcommand{\thepartno}{\alph{partno}}
\renewcommand{\partlabel}{(\thepartno)}

% Fillable table cell: blank in the student copy, the expected value in red on
% the KEY. Fixed-width centered columns keep the blank cells writable.
\newcolumntype{K}[1]{>{\centering\arraybackslash}p{#1}}
\newcommand{\kv}[1]{\ifkey{\color{red}#1}\fi\rule[-0.10in]{0pt}{0.34in}}

% print the answers if this is the key
\ifkey
	\printanswers % if this IS the key
\else
	\noprintanswers  %if this is NOT the key
\fi

% print the header and footer
% need to print key on top if this is the key
\ifkey
	\firstpageheader{}{\color{red}{ KEY}}{}
	\runningheader{\thisclass\hspace{1pt} -- \thisterm }{L\labLesson\ Lab\hspace{1pt}\color{red}{ KEY}}{\thepage\hspace{1pt} of\hspace{1pt} \numpages}
	\firstpagefooter{}{}{}
	\runningfooter{}{}{}
\else
	\firstpageheader{}{}{}
	\runningheader{\thisclass \hspace{1pt} -- \thisterm }{L\labLesson\ Lab\hspace{1pt}}{\thepage\hspace{1pt} of\hspace{1pt} \numpages}
	\firstpagefooter{}{}{}
	\runningfooter{}{}{}
\fi

\runningheadrule


\begin{document}
\pagestyle{headandfoot}   % exam class

%%%%%%%%%%%%%%%%%%%%%%%%%%%%%%%%%%%%%%%%%%%%%%%%%%%%%%
% title block
%%%%%%%%%%%%%%%%%%%%%%%%%%%%%%%%%%%%%%%%%%%%%%%%%%%%%%

\begin{center}
USAF ACADEMY \\
DEPARTMENT OF ELECTRICAL AND COMPUTER ENGINEERING \\
LESSON \labLesson\ LAB - AM and FM Receivers \\
\textbf{Lab Packet} \\
\thisterm \\[4mm]
\end{center}

\vspace{-4pt}
\textbf{Name:} \rule[-2pt]{2.6in}{0.4pt} \hfill
\textbf{Section:} \rule[-2pt]{0.9in}{0.4pt} \hfill
\textbf{Date:} \rule[-2pt]{1.1in}{0.4pt}

\vspace{4pt}

\fullwidth{\textbf{Learning Objective 3.3: } I can implement an FM receiver in GNU
Radio, demodulate a live broadcast FM signal, and decode its RDS data.
\\[2pt]\textbf{Learning Objective 3.4: } I can explain and demonstrate the
trade-off between bandwidth and noise performance for AM vs.\ FM.}

\vspace{2pt}

This packet is the L\labLesson\ Lab, in three parts, and you turn it in once, complete,
before Lesson 8. Do \textbf{Part 1} on your own before the lab period: predict the
receiver settings and the numbers you will measure. \textbf{Part 2} is your own
measurements, on your own laptop and RTL-SDR, recorded in class. Do \textbf{Part 3} on
your own after the lab: compare the measurements with the predictions and answer the
analysis questions.

\textbf{Grading.} One submission, graded against Objectives 3.3 and 3.4 on the
course's 1 / 0 scale, with revision allowed under course policy. Paste screenshots
where a question asks for one.

\textbf{Receive only.} Every flowgraph in this lab receives or simulates. The RTL-SDR
cannot transmit, and nothing in this lab needs a license. Start every receiver with
the volume low.

\textbf{Files.} Download the six L7 flowgraphs from the L7 lesson page into one
folder, with your \texttt{fm\_capture.cfile} from the L5 Lab. Before the period, run
\texttt{L07-fm-test-tx.grc} and \texttt{L07-am-test-tx.grc} once each; they write
\texttt{L07-fm-test.cfile} and \texttt{L07-am-test.cfile} into that folder and stop
on their own. Open \texttt{L07-rds-rx.grc} and confirm that GRC shows no missing
blocks. If it reports the RDS blocks as missing, install gr-rds from the radioconda
prompt with \texttt{conda install gnuradio-rds} and open it again.

\begin{questions}

\fullwidth{\large\textbf{Part 1: Pre-Lab} \normalsize(individual, before the lab period)}

%%%%%%%%%%%%%%%%%%%%%%%%%%%%%%%%%%%%%%%%%%%%%%%%%%%
\question \textbf{Predictions.} You will check each of these at the bench.

\begin{parts}

	\part Fill in the receiver settings for the test file and for your own L5
	capture. The shift is the Frequency Xlating FIR Filter's center frequency.

	\begin{center}
	\renewcommand{\arraystretch}{1.25}
	\begin{tabular}{| l | K{1.2in} | K{1.2in} | K{1.2in} |}
	\hline
	Recording & SDR center (MHz) & Station (MHz) & Shift (kHz) \\\hline
	\texttt{L07-fm-test.cfile} & $100.0$ & $100.3$ & \kv{$+300$} \\\hline
	\texttt{fm\_capture.cfile} & & & \\\hline
	\end{tabular}
	\end{center}

		\begin{solution}[0.5in]
		Shift = station $-$ center. For the L5 capture, answers vary; the shift is
		usually 100 to 400 kHz, because the L5 Lab centered the station slightly off the
		DC spike.
		\end{solution}

	\begin{minipage}[t]{0.58\linewidth}
		\part The FM receiver decimates 2.4~MS/s by 10 and then by 5. Find the two
		rates, and the Quadrature Demod gain that makes 75~kHz of deviation an output
		of 1.
		\begin{solution}[0.9in]
		240 kS/s, then 48 kS/s.
		\eq{ \dfrac{240\,000}{2\pi \times 75\,000} = 0.509 }
		\end{solution}
	\end{minipage}\hfill%
	\begin{minipage}[t]{0.37\linewidth}
		\raggedleft \vspace{12pt}
		\ansbox[1.3in]{240 and 48 kS/s}

		\vspace{4pt}
		gain = \ansbox[1.0in]{$0.509$}
	\end{minipage}

	\begin{minipage}[t]{0.58\linewidth}
		\part The test station's audio is a 1~kHz tone at amplitude 0.6 plus a 2.5~kHz
		tone at 0.4, with pre-emphasis at the transmitter and de-emphasis at the
		receiver. Predict how far below the 1~kHz line the 2.5~kHz line sits on the
		receiver's audio sink.
		\begin{solution}[1.0in]
		De-emphasis undoes pre-emphasis, so the levels are the message's:
		\eq{ 20\mylog{\dfrac{0.4}{0.6}} = -3.5\db }
		\end{solution}
	\end{minipage}\hfill%
	\begin{minipage}[t]{0.37\linewidth}
		\raggedleft \vspace{12pt}
		\ansbox[1.0in]{$3.5\db$ below}
	\end{minipage}

	\begin{minipage}[t]{0.58\linewidth}
		\part RDS sends 1187.5 bits per second in 104-bit groups, and the station name
		takes four groups. Find the shortest time from a clean start to a displayed
		name.
		\begin{solution}[0.8in]
		\eq{ 4 \times \dfrac{104}{1187.5} = 4 \times 87.6 = 350\ \text{ms} }
		\end{solution}
	\end{minipage}\hfill%
	\begin{minipage}[t]{0.37\linewidth}
		\raggedleft \vspace{12pt}
		\ansbox[1.0in]{350 ms}
	\end{minipage}

	\part Find the ATIS frequency of an airport near you in the FAA Chart Supplement
	or on a sectional chart. Choose an SDR center 250~kHz below it, and fill in the AM
	receiver's settings.

	\begin{center}
	\renewcommand{\arraystretch}{1.25}
	\begin{tabular}{| l | K{1.0in} | K{1.0in} | K{1.0in} | K{1.0in} |}
	\hline
	Airport & ATIS (MHz) & SDR center (MHz) & Shift (kHz) & Decimation \\\hline
	& & & \kv{$+250$} & \kv{$50$} \\\hline
	\end{tabular}
	\end{center}

		\begin{solution}[0.4in]
		The decimation is $2\,400\,000 / 48\,000 = 50$ for any station.
		\end{solution}

	\begin{minipage}[t]{0.58\linewidth}
		\part In \texttt{L07-am-vs-fm.grc} the FM channel is $2(25 + 5) = 60$~kHz wide
		and the AM channel 10~kHz. FM reaches threshold at about 10~dB C/N in its own
		bandwidth. Predict the C/N in 10~kHz at which the 25~kHz FM signal reaches
		threshold.
		\begin{solution}[0.9in]
		\eq{ 10 + 10\mylog{\dfrac{60}{10}} = 10 + 7.8 = 17.8\db }
		\end{solution}
	\end{minipage}\hfill%
	\begin{minipage}[t]{0.37\linewidth}
		\raggedleft \vspace{12pt}
		\ansbox[1.0in]{$17.8\db$}
	\end{minipage}

\end{parts}

%%%%%%%%%%%%%%%%%%%%%%%%%%%%%%%%%%%%%%%%%%%%%%%%%%%
\newpage
\fullwidth{\large\textbf{Part 2: Measurements} \normalsize(your own laptop and RTL-SDR, in class)}

\question \textbf{The test file.} Open and run \texttt{L07-wbfm-rx.grc} with its File
Source on \texttt{L07-fm-test.cfile}.

\begin{parts}

	\part Read the audio sink. Record the frequency of each line and how far the
	2.5~kHz line sits below the 1~kHz line.
		\begin{solution}[0.6in]
		Lines at 1 and 2.5 kHz, about 3.5 to 4 dB apart, matching Question 1(c), and
		nothing else above the floor.
		\end{solution}

	\part Drag Station to 99.5~MHz. Describe what you hear, and explain why the station
	at 100.3~MHz disappears.
		\begin{solution}[0.9in]
		A single 600 Hz tone. The Frequency Xlating FIR Filter now shifts 99.5 MHz to
		0 Hz, and the station at 100.3 MHz lands 800 kHz away, far outside the ±90 kHz
		passband.
		\end{solution}

\end{parts}

%%%%%%%%%%%%%%%%%%%%%%%%%%%%%%%%%%%%%%%%%%%%%%%%%%%
\question \textbf{Your L5 capture.} Set the File Source to \texttt{fm\_capture.cfile},
SDR center to the frequency you recorded it at, and Station to the station you
centered.

\begin{parts}

	\part Paste a screenshot of the demodulated multiplex. Mark each part of the
	broadcast multiplex you can see: mono audio, the 19~kHz pilot, the stereo band
	from 23 to 53~kHz, and RDS at 57~kHz.
		\begin{solution}[1.8in]
		Answers vary. A stereo station shows a sharp pilot line at 19 kHz and the stereo
		band; most stations also show the RDS hump at 57 kHz. A mono station shows only
		the audio and possibly RDS.
		\end{solution}

	\part Describe the audio. Your capture is 2~s long, and the File Source repeats it.
		\begin{solution}[0.6in]
		The station's program, repeating every 2 s, in mono.
		\end{solution}

\end{parts}

%%%%%%%%%%%%%%%%%%%%%%%%%%%%%%%%%%%%%%%%%%%%%%%%%%%
\newpage
\question \textbf{Live FM and RDS.} Open \texttt{L07-rds-rx.grc}. Disable the File
Source, enable the RTL-SDR Source, and set SDR center about 300~kHz below a strong
local station and Station to the station. Set RF gain from your L5 Lab gain table.

\begin{parts}

	\part Record what the RDS panel shows. The panel prints the PI code under
	``Program Type'' and the program type under ``PI''; record each under its correct
	name.

	\begin{center}
	\renewcommand{\arraystretch}{1.25}
	\begin{tabular}{| l | p{3.6in} |}
	\hline
	Station (MHz) & \\\hline
	Station name (PS) & \\\hline
	PI code (hex) & \\\hline
	Program type & \\\hline
	Radiotext & \\\hline
	\end{tabular}
	\end{center}

		\begin{solution}[0.3in]
		Answers vary by station.
		\end{solution}

	\part Paste a screenshot of the RDS constellation and the panel. Explain what the
	two clusters show about the receiver.
		\begin{solution}[1.6in]
		Two tight clusters at $\pm 1$ show BPSK with its carrier recovered by the
		Constellation Receiver and its bit timing found by Symbol Sync. A ring or a smear
		would mean the carrier loop has not locked or the signal is too weak.
		\end{solution}

	\begin{minipage}[t]{0.58\linewidth}
		\part Restart the flowgraph and time how long the station name takes to appear.
		Compare with Question 1(d).
		\begin{solution}[0.8in]
		Usually 1 to 3 s: longer than 350 ms, because a station interleaves other group
		types between the name's four groups, and the decoder first has to find the group
		boundaries.
		\end{solution}
	\end{minipage}\hfill%
	\begin{minipage}[t]{0.37\linewidth}
		\raggedleft \vspace{12pt}
		\ansbox[1.0in]{}
	\end{minipage}

\end{parts}

%%%%%%%%%%%%%%%%%%%%%%%%%%%%%%%%%%%%%%%%%%%%%%%%%%%
\question \textbf{ATIS.} Open \texttt{L07-am-rx.grc}. Check it first on
\texttt{L07-am-test.cfile}, then disable the File Source, enable the RTL-SDR Source,
and enter the settings of Question 1(e).

\begin{parts}

	\part Paste a screenshot of the receiver on the live ATIS, and record the
	information letter and the altimeter setting it broadcasts.
		\begin{solution}[1.8in]
		Answers vary. The channel sink shows a carrier line at 0 Hz with voice sidebands
		within about $\pm 3$ kHz, and the audio sink shows speech.
		\end{solution}

	\part Lower RF gain until the voice is hard to understand, and describe how the AM
	audio degrades. You will compare it with FM in Question 6(b).
		\begin{solution}[1.0in]
		The AM voice sinks gradually into a steady hiss and stays intelligible well
		after the hiss is obvious. Answers vary with the station's strength.
		\end{solution}

\end{parts}

%%%%%%%%%%%%%%%%%%%%%%%%%%%%%%%%%%%%%%%%%%%%%%%%%%%
\newpage
\question \textbf{AM and FM in the same noise.} Open and run
\texttt{L07-am-vs-fm.grc}. It needs no hardware.

\begin{parts}

	\part Read the two output SNR meters at each C/N, with the FM deviation at 25~kHz
	and then at 5~kHz. Near threshold the meters wander; record the middle of the
	range you see.

	\begin{center}
	\renewcommand{\arraystretch}{1.25}
	\begin{tabular}{| l | K{0.7in} | K{0.7in} | K{0.7in} | K{0.7in} | K{0.7in} |}
	\hline
	C/N in 10 kHz (dB) & 30 & 25 & 20 & 15 & 10 \\\hline
	AM & \kv{$29.5$} & \kv{$24.5$} & \kv{$19.4$} & \kv{$14.4$} & \kv{$9.2$} \\\hline
	FM, 25 kHz & \kv{$48.5$} & \kv{$44.1$} & \kv{$39.2$} & \kv{$26.1$} & \kv{$9.3$} \\\hline
	FM, 5 kHz & \kv{$35.5$} & \kv{$30.7$} & \kv{$25.5$} & \kv{$20.5$} & \kv{$14.2$} \\\hline
	\end{tabular}
	\end{center}

		\begin{solution}[0.5in]
		Within about 1 dB of these values; the FM entries near threshold vary the most.
		\end{solution}

	\part Set the deviation to 25~kHz and Listen To to FM. Lower C/N from 30 to 10~dB
	and describe the sound at 20, 15, and 10~dB.
		\begin{solution}[1.0in]
		Clean at 20 dB. Clicks begin near 15 dB and become frequent at 10 dB, where the
		tone sounds rough. The clicks are the FM threshold. AM, in Question 5(b), only
		grew steadily noisier.
		\end{solution}

\end{parts}

%%%%%%%%%%%%%%%%%%%%%%%%%%%%%%%%%%%%%%%%%%%%%%%%%%%
\newpage
\fullwidth{\large\textbf{Part 3: Writeup} \normalsize(individual, after the lab)}

\question \textbf{Analysis.}

\begin{parts}

	\part Your mono receiver plays a stereo station correctly. Using your multiplex
	screenshot from Question 3(a), explain why.
		\begin{solution}[1.1in]
		The mono audio, the sum of the left and right channels, fills 0 to 15 kHz by
		itself. The pilot, the stereo difference signal, and RDS all sit above 15 kHz, so
		the receiver's 15 kHz low-pass filter removes them and plays the sum.
		\end{solution}

	\part Compare where your 25~kHz FM measurements lost their advantage over AM with
	your prediction in Question 1(f). Then use your table to explain why the 5~kHz
	signal beats the 25~kHz signal at 10~dB C/N but loses at 30~dB.
		\begin{solution}[1.6in]
		The 25 kHz signal leads AM by about 19 to 20 dB down to 20 dB C/N, has lost
		about 7 dB of that lead at 15 dB, and is no better than AM at 10 dB, consistent
		with a threshold near the predicted 17.8 dB. Wider deviation gives a larger
		output SNR above threshold, 48.5 against 35.5 dB at 30 dB C/N. It also needs a
		wider channel, 60 kHz against 20 kHz, which admits 4.8 dB more noise, so the
		25 kHz signal falls below threshold first. The 5 kHz signal is still above its
		threshold at 10 dB and keeps its 5 dB lead over AM.
		\end{solution}

	\part Aircraft use AM, and broadcasters use wideband FM. Using your measurements
	from Questions 5 and 6, explain one reason for each choice.
		\begin{solution}[1.2in]
		Broadcast FM: a broadcast station is strong and has a 200 kHz channel, so it
		stays far above threshold and the wide deviation gives a much cleaner output than
		AM. Aircraft AM: a distant aircraft is weak, and an AM signal degrades gradually
		and stays intelligible where FM would break into clicks. Two aircraft transmitting
		at once are also both heard on AM.
		\end{solution}

\end{parts}

\end{questions}

\end{document}
